\chapter{Magnesium (Blood)}

\section*{What Is This Test?}
The magnesium blood test measures the concentration of magnesium in the blood serum or plasma. Magnesium is the fourth most abundant cation in the human body and the second most abundant intracellular cation after potassium. Approximately 50--60\% of total body magnesium (22--26 grams) is stored in bone as hydroxyapatite-adsorbed surface ions, 25--30\% in muscle, and less than 1\% in the extracellular fluid, making serum magnesium a poor indicator of total body stores. In blood, magnesium exists in three forms: (1) ionized magnesium (free, physiologically active, approximately 55--70\%), (2) protein-bound magnesium (approximately 20--30\%, primarily bound to albumin), and (3) magnesium complexed with anions such as phosphate, citrate, and bicarbonate (5--15\%). Magnesium is an essential cofactor for over 300 enzymatic reactions, including all ATP-requiring processes (ATP must be complexed with Mg\textsuperscript{2+} as Mg-ATP to be biologically active), DNA and protein synthesis, oxidative phosphorylation, glycolysis, and neuromuscular transmission \cite{debaaij2015magnesium}. Magnesium regulates calcium flux across cell membranes by competing with calcium for membrane binding sites and modulating calcium channels; it is often called ``nature's calcium channel blocker.'' Magnesium is critical for parathyroid hormone (PTH) secretion---hypomagnesemia impairs PTH secretion and causes end-organ resistance to PTH, resulting in functional hypoparathyroidism with hypocalcemia that is refractory to calcium replacement until magnesium is corrected. Renal magnesium handling is the primary regulator of magnesium balance: approximately 70--80\% of filtered magnesium is reabsorbed in the thick ascending limb of Henle via paracellular transport driven by the positive luminal potential generated by the Na-K-2Cl cotransporter (the target of loop diuretics), and 5--10\% in the distal convoluted tubule via the TRPM6 magnesium channel. Magnesium is absorbed in the small intestine (30--50\% of dietary intake, primarily in the ileum) by both paracellular and active (TRPM6) transport.

\section*{Alternative Names}
Mg; Magnesium, Serum; Serum Magnesium; Plasma Magnesium; Mg2+; Ionized Magnesium; RBC Magnesium (erythrocyte magnesium for tissue stores)

\section*{Principle}
Serum magnesium is measured by spectrophotometric methods using colorimetric complex-forming reagents. The most common methods are: (1) Calmagite (1-(1-hydroxy-4-methyl-2-phenylazo)-2-naphthol-4-sulfonic acid) method: magnesium reacts with calmagite in alkaline solution to form a reddish-violet complex absorbing at 530--540 nm. Calcium interference is prevented by adding EGTA (egtazic acid) or strontium salts. (2) Xylidyl blue (Magonsulfonate) method: magnesium forms a complex with the sulfonated dye in alkaline buffer producing color measured at 600 nm. (3) Methylthymol blue method: magnesium reacts with methylthymol blue in alkaline solution to form a blue complex absorbing at 600 nm. (4) Enzymatic methods using hexokinase (which requires Mg-ATP as substrate) or glycerol kinase, where the rate of the enzymatic reaction is proportional to magnesium concentration. Ionized magnesium can be measured by ion-selective electrode (ISE) with ETH 7025 or related neutral ionophores; however, ionized magnesium ISEs are less widely available and less reliable than calcium ISEs, and are primarily used in research settings. Atomic absorption spectrophotometry was the historical gold standard but is rarely used clinically. The test is performed on serum separator tubes using spectrophotometric methods, not by hematology analyzers \cite{tietz2023textbook}.

\section*{History}
Magnesium was first isolated by Humphry Davy in 1808. The essential role of magnesium in biology was recognized in the early 20th century when McCollum and colleagues demonstrated magnesium deficiency caused vasodilation and tetany in animals (1930s). The first clinical recognition of magnesium deficiency in humans was by Hirschfelder in 1934, who documented low serum magnesium in patients with various disorders. The calmagite colorimetric method was developed in the 1960s. The magnesium ion-selective electrode was introduced in the 1980s. The discovery of TRPM6 as the distal tubule and intestinal magnesium channel was published in 2002, explaining the pathogenesis of hereditary hypomagnesemia with secondary hypocalcemia.

\section*{Why This Test Is Ordered}
\begin{itemize}
  \item Evaluation of neuromuscular irritability: tremors, fasciculations, muscle cramps, hyperreflexia, positive Chvostek and Trousseau signs (similar to hypocalcemia), tetany, seizures
  \item Evaluation of cardiac arrhythmias: atrial fibrillation, ventricular ectopy, torsades de pointes (polymorphic VT), prolonged QT interval. Hypomagnesemia is present in 20--35\% of ICU patients with cardiac arrhythmias
  \item Refractory hypokalemia: hypomagnesemia causes renal potassium wasting that is resistant to potassium replacement until magnesium is corrected. Always check magnesium in unexplained or refractory hypokalemia (Mg <2.0 mg/dL or <1.7 mg/dL causes significant renal potassium loss)
  \item Unexplained hypocalcemia: hypomagnesemia impairs PTH secretion and causes PTH resistance; hypocalcemia does not respond to calcium replacement alone
  \item Monitoring of patients on medications causing magnesium depletion: loop diuretics, thiazide diuretics, proton pump inhibitors (PPIs, especially with prolonged use >1 year), aminoglycosides, amphotericin B, cisplatin, carboplatin, cyclosporine, tacrolimus, cetuximab, panitumumab, pentamidine, foscarnet
  \item Gastrointestinal disorders causing magnesium malabsorption or loss: chronic diarrhea, Crohn's disease, celiac disease, short bowel syndrome, intestinal resection, chronic pancreatitis, steatorrhea, prolonged nasogastric suction, acute pancreatitis
  \item Alcohol use disorder: magnesium deficiency in 30--60\% of chronic alcoholics due to poor dietary intake, GI losses (diarrhea, vomiting), and renal magnesium wasting from alcohol's direct tubular effect
  \item Preeclampsia and eclampsia: magnesium sulfate is the standard treatment for seizure prophylaxis in severe preeclampsia and eclampsia; serum magnesium levels are monitored to maintain therapeutic range (4--7 mg/dL)
  \item Evaluation of suspected Gitelman syndrome (hypokalemia, hypomagnesemia, metabolic alkalosis, hypocalciuria---genetic defect in distal Na-Cl cotransporter) and Bartter syndrome
  \item Diabetes: hypomagnesemia is common due to renal magnesium wasting from osmotic diuresis; magnesium deficiency worsens insulin resistance
  \item Refractory ventricular arrhythmias/cardiac arrest: IV magnesium is first-line for torsades de pointes (2 g MgSO4 IV push over 1--2 minutes)
\end{itemize}

\section*{Primary vs Secondary Test}
Serum magnesium is a secondary test. It is not part of the basic or comprehensive metabolic panel and must be specifically ordered. It should be checked in patients with unexplained hypokalemia, hypocalcemia, cardiac arrhythmias, chronic diarrhea, alcoholism, or those on medications causing magnesium depletion \cite{tierney2023current}.

\section*{How This Test Is Performed}
A healthcare professional draws blood from a vein into a serum separator tube (gold-top) or lithium heparin tube (green-top). Prolonged tourniquet use should be minimal. EDTA tubes (purple-top) must NOT be used because EDTA chelates magnesium. The sample is centrifuged and analyzed on automated chemistry analyzers using colorimetric methods. Results are usually available within 1--2 hours. Unlike calcium, magnesium measurement does not require anaerobic handling. The blood draw is a standard venipuncture taking less than 5 minutes.

\section*{What Preparation Is Needed}
No fasting is required. Inform the healthcare provider of all medications, especially diuretics (loop and thiazide), proton pump inhibitors (omeprazole, esomeprazole, pantoprazole---associated with hypomagnesemia with prolonged use >1 year), aminoglycosides, cisplatin, cyclosporine, tacrolimus, cetuximab, and magnesium-containing supplements or antacids (which can elevate measured magnesium if taken shortly before the test). Avoid magnesium-containing laxatives or antacids on the morning of the test. Prolonged tourniquet time should be avoided as it may cause slight local tissue release of magnesium. Samples should be processed within 4 hours at room temperature or 24 hours refrigerated. Hemolysis must be avoided because erythrocyte magnesium concentration is approximately 3 times higher than serum magnesium.

\section*{Contraindications and Precautions}
No absolute contraindications. Critical pre-analytical considerations: (1) Hemolysis: red blood cells contain magnesium at approximately 5.0--6.5 mg/dL (3$\times$ serum concentration); even mild hemolysis can significantly elevate serum magnesium and cause pseudohyper- or pseudonormomagnesemia in a truly deficient patient. Hemolyzed samples must be rejected and redrawn \cite{workinger2018challenges}. (2) EDTA tube: EDTA chelates magnesium causing pseudohypomagnesemia. (3) Delayed centrifugation: magnesium is slowly released from red blood cells over time. (4) Hypoalbuminemia: total serum magnesium is partially albumin-bound, so low albumin can decrease total magnesium without affecting ionized magnesium. Unlike calcium, routine albumin correction is not standard practice for magnesium. (5) Hemoconcentration from prolonged tourniquet use can slightly elevate magnesium. (6) Medications: loop diuretics and thiazides deplete magnesium (block renal reabsorption); PPIs with prolonged use (>1 year) can impair intestinal magnesium absorption through reduced TRPM6 activity; cisplatin causes direct renal tubular magnesium wasting; aminoglycosides block calcium-sensing receptor-mediated magnesium reabsorption in thick ascending limb; cetuximab and panitumumab (EGFR inhibitors) block TRPM6 expression causing severe hypomagnesemia. (7) Contrast dye (gadolinium) can interfere with some colorimetric methods \cite{wallach2022interpretation}.

\section*{Risks}
Minimal risk from venipuncture. Mild pain or bruising at the site resolves in 1--2 days. Rare: hematoma, infection, nerve injury, vasovagal syncope. No additional risk specific to magnesium measurement.

\section*{What Happens After the Test}
Results are typically available in 1--2 hours. Critical values: magnesium <1.0 mg/dL (risk of seizures, tetany, refractory ventricular arrhythmias) or >5.0 mg/dL (risk of respiratory depression, loss of deep tendon reflexes, cardiac arrest). In preeclampsia/eclampsia patients receiving IV magnesium sulfate, therapeutic range is 4--7 mg/dL; loss of deep tendon reflexes occurs at 8--12 mg/dL (first sign of toxicity), respiratory depression at 12--15 mg/dL, and cardiac arrest at >15--18 mg/dL. When hypomagnesemia is found, oral replacement is limited by GI tolerance (diarrhea from unabsorbed magnesium); IV magnesium sulfate is often required for severe deficiency. Hypomagnesemia should always prompt checking of potassium and calcium simultaneously, as correction of magnesium deficiency is required for effective potassium and calcium replacement.

\section*{Understanding Results}

\subsection*{Below Normal}
\begin{itemize}
  \item \textbf{Hypomagnesemia (Mg <1.7 mg/dL or <0.7 mmol/L):} Present in approximately 10--20\% of hospitalized patients and up to 60\% of ICU patients. Mild 1.2--1.7 mg/dL; Moderate 0.5--1.2 mg/dL; Severe <0.5 mg/dL \cite{agus1999hypomagnesemia}.
  \item \textbf{Gastrointestinal losses (most common broad category):} Chronic diarrhea (any cause---inflammatory bowel disease, infectious, osmotic, secretory), malabsorption syndromes (celiac disease, Crohn's disease, short bowel syndrome, intestinal resection), steatorrhea (magnesium forms soaps with fatty acids), acute and chronic pancreatitis, prolonged nasogastric suction, biliary or pancreatic fistula
  \item \textbf{Renal losses:} Loop diuretics (furosemide, bumetanide, torsemide---block Na-K-2Cl cotransporter in thick ascending limb, collapsing the positive luminal potential that drives paracellular magnesium reabsorption); thiazide diuretics (reduce TRPM6 expression in distal tubule); osmotic diuresis (hyperglycemia, mannitol); medications (cisplatin---causes severe, often permanent hypomagnesemia; carboplatin; aminoglycosides; amphotericin B; cyclosporine and tacrolimus---reduce TRPM6 expression; cetuximab and panitumumab---EGFR inhibitors block TRPM6 trafficking; pentamidine; foscarnet); alcohol (direct tubular toxicity); hypercalcemia (calcium competes with magnesium for reabsorption in thick ascending limb); recovery phase of acute tubular necrosis; post-obstructive diuresis; Gitelman syndrome (mutation in SLC12A3 encoding distal Na-Cl cotransporter---hypokalemic metabolic alkalosis, hypomagnesemia, hypocalciuria); Bartter syndrome (mutation in thick ascending limb transporters)
  \item \textbf{Proton pump inhibitor (PPI)-associated hypomagnesemia:} PPIs (omeprazole, esomeprazole, pantoprazole, lansoprazole) with prolonged use (>1 year, especially in elderly) inhibit intestinal magnesium absorption via TRPM6; can cause severe, symptomatic hypomagnesemia refractory to oral replacement \cite{fda2011low}. FDA issued a safety warning in 2011. Mechanism: reduced intestinal pH impairs TRPM6-mediated magnesium absorption
  \item \textbf{Inadequate intake:} Malnutrition, alcoholism (poor diet + direct renal magnesium wasting from alcohol), prolonged TPN without adequate magnesium supplementation, refeeding syndrome (refeeding stimulates cellular magnesium uptake along with glucose, phosphate, potassium)
  \item \textbf{Endocrine disorders:} Diabetes mellitus (osmotic diuresis from hyperglycemia causes renal magnesium wasting---hypomagnesemia in 25--39\% of poorly controlled diabetics; magnesium deficiency worsens insulin resistance); primary hyperparathyroidism (increased urinary magnesium excretion); hungry bone syndrome after parathyroidectomy
  \item \textbf{Clinical manifestations:} Neuromuscular: hyperreflexia, positive Chvostek/Trousseau signs, muscle cramps, fasciculations, tetany, seizures, vertical nystagmus. Cardiac: prolonged QT, PR, and QRS intervals; atrial fibrillation; ventricular ectopy; torsades de pointes. Metabolic: hypokalemia (refractory to potassium replacement), hypocalcemia (impaired PTH secretion and end-organ resistance). Magnesium <1.2 mg/dL: severe symptoms with high risk of arrhythmia.
\end{itemize}

\subsection*{Above Normal}
\begin{itemize}
  \item \textbf{Hypermagnesemia (Mg >2.2 mg/dL or >0.9 mmol/L):} Uncommon in patients with normal renal function because the kidney can excrete large magnesium loads. Almost exclusively seen in patients with renal impairment plus excess magnesium intake.
  \item \textbf{Renal failure (most common cause):} Acute kidney injury or chronic kidney disease with GFR <30 mL/min impairs magnesium excretion. In advanced CKD, hypermagnesemia develops from dietary magnesium plus medications. Dialysis patients are at risk if dialysate magnesium concentration is too high. GI magnesium absorption continues even with negligible renal excretion.
  \item \textbf{Excessive intake in setting of renal impairment:} Magnesium-containing antacids (magnesium hydroxide, magnesium carbonate), magnesium-containing laxatives (magnesium citrate---2.5 g magnesium per dose; milk of magnesia), magnesium-containing enemas, Epsom salt (magnesium sulfate) use. Even small therapeutic doses can cause severe hypermagnesemia in patients with ESRD.
  \item \textbf{Iatrogenic administration:} IV magnesium sulfate for preeclampsia/eclampsia (therapeutic range 4--7 mg/dL, toxicity at >8 mg/dL with loss of deep tendon reflexes), TPN with excessive magnesium, magnesium-containing urologic irrigating solutions. Neonatal hypermagnesemia occurs in infants born to mothers receiving IV MgSO4 for preeclampsia.
  \item \textbf{Other:} Lithium therapy (lithium can cause mild hypermagnesemia by competing with magnesium for renal handling); hypothyroidism; Addison's disease; milk-alkali syndrome; familial hypocalciuric hypercalcemia (FHH)
  \item \textbf{Clinical manifestations (progressive with rising Mg):} 4--7 mg/dL: Therapeutic range (preeclampsia). 5--8 mg/dL: Nausea, flushing, headache, mild hypotension. 8--12 mg/dL: Loss of deep tendon reflexes (first sign of magnesium toxicity), somnolence, bradycardia, hypotension, prolonged PR and QRS intervals. 12--15 mg/dL: Respiratory depression/paralysis, coma. >15--18 mg/dL: Cardiac arrest (asystole, complete heart block). The earliest sign of magnesium toxicity is loss of patellar (knee-jerk) deep tendon reflexes.
  \item \textbf{Treatment:} IV calcium gluconate or calcium chloride (direct antagonist of magnesium at the neuromuscular junction and cardiac membrane). For severe hypermagnesemia in renal failure, hemodialysis is definitive therapy. For patients with intact renal function, IV normal saline with loop diuretics (furosemide) enhances renal magnesium excretion.
\end{itemize}

\section*{Normal Results}
\begin{itemize}
  \item \textbf{Adults:} 1.7--2.2 mg/dL (0.7--0.9 mmol/L, 1.4--1.8 mEq/L) \cite{fischbach2023manual}
  \item \textbf{Children (1--16 years):} 1.7--2.1 mg/dL
  \item \textbf{Infants (1--12 months):} 1.5--2.2 mg/dL
  \item \textbf{Newborns (0--7 days):} 1.5--2.3 mg/dL
  \item \textbf{RBC magnesium (tissue stores):} 4.3--6.5 mg/dL (better reflects total body stores than serum Mg)
  \item \textbf{Ionized magnesium:} 0.5--0.7 mmol/L (1.2--1.7 mg/dL)
  \item \textbf{Critical low value:} <1.0 mg/dL
  \item \textbf{Critical high value:} >5.0 mg/dL (non-obstetric); >8.0 mg/dL (obstetric/IV MgSO4)
  \item \textbf{Therapeutic range (MgSO4 for eclampsia):} 4--7 mg/dL
  \item \textbf{Conversion:} 1 mg/dL = 0.411 mmol/L; 1 mmol/L = 2.43 mg/dL
\end{itemize}

\section*{Follow-up Tests}
\begin{itemize}
  \item \textbf{Potassium level:} Always check simultaneously---hypomagnesemia causes renal potassium wasting; potassium replacement is ineffective until magnesium is corrected to >2.0 mg/dL.
  \item \textbf{Calcium level:} Hypomagnesemia impairs PTH secretion and causes end-organ PTH resistance; hypocalcemia is refractory to calcium replacement until magnesium is corrected.
  \item \textbf{ECG:} Assess QT interval, look for U waves, ventricular ectopy, atrial fibrillation. Hypomagnesemia prolongs QT (risk of torsades de pointes). Hypermagnesemia prolongs PR and QRS at toxic levels.
  \item \textbf{Renal function (BUN, creatinine):} Essential because hypermagnesemia rarely occurs with normal renal function; CKD is the most common predisposition.
  \item \textbf{24-hour urine magnesium:} Distinguishes renal from GI magnesium losses. In hypomagnesemia, 24-hour urine Mg >24 mg (or fractional excretion >2\%) indicates inappropriate renal magnesium wasting (loop diuretics, cisplatin, Gitelman syndrome). Urine Mg <12 mg/24h indicates GI loss or inadequate intake with appropriate renal conservation.
  \item \textbf{Plasma renin activity and aldosterone:} For suspected Gitelman or Bartter syndrome (hypokalemia, metabolic alkalosis, hypomagnesemia). Bartter: hypercalciuria; Gitelman: hypocalciuria.
  \item \textbf{RBC magnesium:} More accurately reflects total body magnesium stores than serum magnesium, because magnesium is primarily intracellular. May identify magnesium deficiency in patients with normal serum magnesium but high clinical suspicion.
  \item \textbf{PTH level:} If hypomagnesemia with hypocalcemia; determines if hypoparathyroidism is present.
  \item \textbf{25-hydroxyvitamin D:} If malabsorption is suspected as cause of magnesium deficiency.
  \item \textbf{Deep tendon reflex assessment:} Essential bedside monitoring for magnesium toxicity in patients receiving IV MgSO4; loss of patellar reflexes is earliest sign (Mg >8 mg/dL).
\end{itemize}

\section*{References}
\bibliographystyle{unsrtnat}
\bibliography{references}

\clearpage
